\section*{Introduzione}
%--------1---------2---------3---------4---------5---------6---------7---------8
Nel mondo mainframe una elaborazione batch normalmente è divisa in più parti
(step), e questi sono coordinati da uno script di controllo (JCL).
%--------1---------2---------3---------4---------5---------6---------7---------8
Questa libreria mette a disposizione delle classi per facilitare lo sviluppo
delle elaborazioni dei singoli step e di coordinale l'\,esecuzione degli step.
%--------1---------2---------3---------4---------5---------6---------7---------8

%--------1---------2---------3---------4---------5---------6---------7---------8
\begin{elisting}[!htb]
\begin{javacode}
public Integer call() {
    return Job.create("job01")
        .execPgm("step01", step01::run)
        .complete();
}
\end{javacode}
\caption{esempio minimale job}
\label{lst:demoJob}
\end{elisting}
%--------1---------2---------3---------4---------5---------6---------7---------8
Un esempio minimale di batch con un singolo step è nel lis.~\ref{lst:demoJob}
dove \texttt{job01} è il nome del job, \texttt{step01} è il come dello step e
\texttt{step01::run} è il riferimento al metodo che esegue lo step.
%--------1---------2---------3---------4---------5---------6---------7---------8

%--------1---------2---------3---------4---------5---------6---------7---------8
L'\,esecuzione del job produce nel log un report finale del tipo del
lis.~\ref{lst:demoLog}.
%--------1---------2---------3---------4---------5---------6---------7---------8
\begin{elisting}[!htb]
\begin{Verbatim}[fontsize=\small,frame=single]
  Name   !  rc  !     Date-Time Start/Skip      !         Date-Time End         !       Lapse
---------+------+-------------------------------+-------------------------------+-------------------
step01...!    0 ! 2024-05-17T14:59:12.03243396  ! 2024-05-17T14:59:12.590157083 ! 00:00:00.557723123
---------+------+-------------------------------+-------------------------------+-------------------
job01....!    0 ! 2024-05-17T14:59:12.028019213 ! 2024-05-17T14:59:12.590998545 ! 00:00:00.562979332
\end{Verbatim}
\caption{esempio report fine esecuzione job}
\label{lst:demoLog}
\end{elisting}

%--------1---------2---------3---------4---------5---------6---------7---------8
La libreria è divisa in 5 moduli:
\begin{itemize}
%--------1---------2---------3---------4---------5---------6---------7---------8
\item \texttt{promethium-batch-fault} contiene le classi per la gestione degli
  errori e i codici di errori;
%--------1---------2---------3---------4---------5---------6---------7---------8
\item \texttt{promethium-batch-job} contiene le classi per la definizione del
  job e la coordinazione degli step (richiede \texttt{promethium-batch-fault});
%--------1---------2---------3---------4---------5---------6---------7---------8
\item \texttt{promethium-batch-pgm-core} contiene le classi usate per la gestione
  dei dati di input e di output (richiede \texttt{promethium-batch-fault});
%--------1---------2---------3---------4---------5---------6---------7---------8
\item \texttt{promethium-batch-pgm-maven-plugin} è un plugin per la generazione
  delle classi per gestire la elaborazione dei dati all'\,interno di uno step
  (il codice generato richiede \texttt{promethium-batch-pgm-core});
%--------1---------2---------3---------4---------5---------6---------7---------8
\item \texttt{promethium-batch-pgm} contiene le classi generati dal plugin
  nel caso di elaborazioni $1 \to N$ con $N$ da $0$ a $12$, e per le
  elaborazioni $N_I \to N_O$ con $N_I$ da $2$ a $3$ e $N_O$ da $0$ a $8$
  (richiede \texttt{promethium-batch-pgm-core}).
\end{itemize}

%--------1---------2---------3---------4---------5---------6---------7---------8
Se non è opportuno sviluppare gli step con la struttura delle classi
\verb!pgm*! si può utilizzare solo la prima libreria
\texttt{promethium-batch-job}.

%--------1---------2---------3---------4---------5---------6---------7---------8
Se l'\,elaborazione usa un singolo step si può ritenere di poca utilità
l'\,introduzione della struttura della libreria \verb!job! e utilizzare soltanto
la libreria \texttt{promethium-batch-pgm}.

%--------1---------2---------3---------4---------5---------6---------7---------8
I valori $N, N_I, N_O$ usati dalla libreria \verb!pgm! sono normalmente
sufficienti, aumentare questi valori porta a un aumento \textsl{quadratico}
delle dimensioni di questa libreria.
%--------1---------2---------3---------4---------5---------6---------7---------8
Se fossero necessari valori $N, N_I, N_O$ superiori a quelli usati dalla
libreria \texttt{promethium-batch-pgm}, è possibile utilizzare al suo posto il
plugin che è in grado di generare le classi \textsl{pgm} in modo personalizzato.
Anziché implementate un range completo (\texttt{0..12}) è possibile implementare
un range ristretto (\texttt{20..24}), o un numero preciso di valori
(\texttt{5,12,18}).
In questo modo vengono generate solo le implementazioni che servono e la
dimensione del codice rimane limitata.
